\documentclass[12pt]{article}
\usepackage{amsmath,amssymb,amsthm}
\usepackage[margin=1in]{geometry}

\newtheorem{theorem}{Theorem}
\newtheorem{lemma}[theorem]{Lemma}
\newtheorem{corollary}[theorem]{Corollary}
\newtheorem{definition}[theorem]{Definition}

\title{Problem 301: Nim}
\author{Project Euler Solution}
\date{}

\begin{document}
\maketitle

\section{Problem Statement}
Nim is a two-player combinatorial game played with heaps of stones under the normal play convention: players alternate turns removing any positive number of stones from a single heap, and the player who removes the last stone wins.

Consider the three-heap Nim position $(n, 2n, 3n)$. How many positive integers $n \le 2^{30}$ satisfy $W(n, 2n, 3n) = 0$ (first player loses)?

By the Sprague--Grundy theorem, a Nim position $(a,b,c)$ is a P-position (losing for the player to move) if and only if $a \oplus b \oplus c = 0$. The problem thus reduces to counting $n \le 2^{30}$ with $n \oplus 2n \oplus 3n = 0$.

\section{Mathematical Development}

\begin{lemma}[Carry-free addition]\label{lem:xor-add}
For non-negative integers $a$ and $b$, $a \oplus b = a + b$ if and only if $a \mathbin{\&} b = 0$.
\end{lemma}

\begin{proof}
Write $a = \sum_i a_i 2^i$ and $b = \sum_i b_i 2^i$ with $a_i, b_i \in \{0,1\}$. Bitwise XOR computes columnwise addition modulo~2 without carries: $(a \oplus b)_i = a_i + b_i \pmod{2}$. The ordinary sum $a + b = \sum_i (a_i + b_i)2^i$ holds without higher-order corrections precisely when no carry is generated, i.e., $a_i + b_i \le 1$ for all~$i$. This is equivalent to $a_i \cdot b_i = 0$ for all~$i$, i.e., $a \mathbin{\&} b = 0$.
\end{proof}

\begin{theorem}[Characterization of losing positions]\label{thm:consec}
$n \oplus 2n \oplus 3n = 0$ if and only if $n$ has no two consecutive $1$-bits in its binary representation.
\end{theorem}

\begin{proof}
Since $3n = n + 2n$, the condition $n \oplus 2n \oplus 3n = 0$ is equivalent to $n \oplus 2n = 3n = n + 2n$. By Lemma~\ref{lem:xor-add}, this holds if and only if $n \mathbin{\&} 2n = 0$.

Multiplication by~$2$ is a left shift by one bit: bit~$i$ of $2n$ equals bit~$(i{-}1)$ of~$n$. Writing $n = \sum_i b_i 2^i$, we obtain
\[
n \mathbin{\&} 2n = \sum_{i \ge 1} b_i \cdot b_{i-1} \cdot 2^i.
\]
This vanishes if and only if $b_i \cdot b_{i-1} = 0$ for all $i \ge 1$, which is precisely the condition that no two consecutive bits of~$n$ are both~$1$.
\end{proof}

\begin{lemma}[Fibonacci counting]\label{lem:fib-count}
Let $a_k$ denote the number of binary strings of length~$k$ (with possible leading zeros) having no two consecutive $1$-bits. Then $a_k = F_{k+2}$, where $F_1 = F_2 = 1$ and $F_n = F_{n-1} + F_{n-2}$.
\end{lemma}

\begin{proof}
A valid $k$-bit string either begins with~$0$ (giving $a_{k-1}$ continuations) or begins with~$10$ (giving $a_{k-2}$ continuations). Hence $a_k = a_{k-1} + a_{k-2}$ for $k \ge 3$. The base cases $a_1 = 2$ and $a_2 = 3$ match $F_3 = 2$ and $F_4 = 3$. By induction, $a_k = F_{k+2}$.
\end{proof}

\begin{corollary}\label{cor:count-range}
The number of non-negative integers in $\{0, 1, \ldots, 2^k - 1\}$ with no two consecutive $1$-bits is $F_{k+2}$.
\end{corollary}

\begin{proof}
The integers in $\{0,\ldots,2^k-1\}$ correspond bijectively to $k$-bit strings. Apply Lemma~\ref{lem:fib-count}.
\end{proof}

\begin{theorem}[Main result]\label{thm:main}
The number of positive integers $n \le 2^{30}$ with no two consecutive $1$-bits is $F_{32} = 2{,}178{,}309$.
\end{theorem}

\begin{proof}
By Corollary~\ref{cor:count-range}, the qualifying integers in $\{0, 1, \ldots, 2^{30}-1\}$ number $F_{32}$. Excluding $n = 0$ gives $F_{32} - 1$. The integer $n = 2^{30} = 1\underbrace{00\cdots0}_{30}$ trivially has no consecutive $1$-bits. Including it yields $F_{32} - 1 + 1 = F_{32} = 2{,}178{,}309$.
\end{proof}

\section{Algorithm}
\begin{verbatim}
function solve():
    a <- 1; b <- 1
    for i = 3 to 32:
        (a, b) <- (b, a + b)
    return b
\end{verbatim}

\section{Complexity Analysis}
\begin{itemize}
    \item \textbf{Time:} $O(1)$ --- compute $F_{32}$ via exactly 30 additions.
    \item \textbf{Space:} $O(1)$ --- two integer variables suffice.
\end{itemize}

\section{Answer}

\[
\boxed{2178309}
\]

\end{document}
